\section{The limiting Bui corner}
\label{sec:bui-corner}

This section is a boundary calculation and is not used in the proof of
\Cref{thm:main}.  Write Bui's two mollifier lengths as
\begin{equation}\label{eq:Bui-box}
 H\leq T^{4/7-\eps_B},\qquad
 K\leq T^{1/4-\eps_B}.
\end{equation}
Then
\begin{equation}\label{eq:Bui-fixed-Weil-exponent}
 \frac32\left(\frac47-\eps_B\right)
 +\left(\frac14-\eps_B\right)
 =\frac{31}{28}-\frac52\eps_B.
\end{equation}
It is less than one precisely when
\begin{equation}\label{eq:Bui-threshold}
                              \eps_B>\frac{3}{70}.
\end{equation}
Thus \Cref{thm:main} already covers every Bui box with this strict margin.

If the margin is suppressed and only the limiting exponents
$(\theta_1,\theta_3)=(4/7,1/4)$ are compared, the nonzero-mode estimate is
\begin{equation}\label{eq:Bui-limiting-budget}
                    KH^{3/2}=T^{31/28+o(1)}.
\end{equation}
To reach an error $T^{1-\delta}$ through the same reduction would require an
additional saving of
\begin{equation}\label{eq:Bui-needed-saving}
                              T^{3/28+\delta}
\end{equation}
in the limiting exponent budget.  The pointwise completed Weil bound in
\Cref{lem:reciprocal-sum} does not provide it for two arbitrary coefficient
sequences.

An averaged reciprocal or Kloosterman-fraction estimate retaining both
coefficient sequences could be a sufficient new input for the nonzero
frequencies, but no such estimate is assumed or proved here.  Nor would it by
itself settle every analytic detail at the corner: the endpoint-deletion and
full-kernel estimates in \Cref{sec:endpoint} were proved and audited under
\eqref{eq:main-range}, and their uniformity at the limiting corner has not
been established.  It is therefore not valid to take a minimum over different
Poisson orientations or to describe a single averaged estimate as the unique
remaining dependency.

Finally, the cubic product in \eqref{eq:I-def} is the ``twisted third
moment'' appearing in Bui's mollifier cross term.  It is not the sixth moment
$\int|\zeta(1/2+it)|^6\dd t$, and \Cref{thm:main} does not imply a sixth-moment
asymptotic.
